\begin{frame}
  \titlepage
\end{frame}

\begin{frame}{Papers}
  \vspace{-0.25em}
  {\fontsize{7.9}{9.1}\selectfont
  \begin{block}{Paper A}
    Kalicanin Dimitrov, M., Dimitrov, M., Malyarenko, A., and Ni, Y. (2025).
    \emph{Almost-Exact Simulation Scheme for Heston-type Models: Bermudan and American Option Pricing}.
    Manuscript submitted for publication.
  \end{block}

  \vspace{-0.1em}
  \begin{block}{Paper B}
    Haida, B., Dimitrov, M. K., Marmaras, T., and Ni, Y. (2026).
    \emph{Almost Exact Mixed Scheme for Gatheral Double Mean Reverting Model}.
    In S. Silvestrov and A. Malyarenko (Eds.),
    \emph{Mathematical Structures I: Stochastic and Analytic Structures with Applications},
    Springer Proceedings in Mathematics and Statistics, Vol. 532. Springer, Cham.
    Accepted for publication.
  \end{block}

  \vspace{-0.1em}
  \begin{block}{Paper C}
    M. Kalicanin Dimitrov and Y. Ni (2026).
    \emph{A Hybrid LSMC--PDE Method for Bermudan Option Pricing under the Gatheral Double Mean-Reverting Model}.
    Submitted to International Journal of Applied and Computational Mathematics (currently under review).
  \end{block}}
\end{frame}

\begin{frame}{One Numerical Idea Across Three Papers}
  \keymessage{One common numerical idea connects the three papers: use the structure of the stochastic-volatility model to improve either path simulation or continuation-value approximation.}

  \vspace{0.6em}
  \begin{columns}[T,onlytextwidth]
    \begin{column}{0.31\textwidth}
      \begin{block}{Paper A}
        \paperlabel{AES}\\
        Heston and double Heston. CIR variance sampling is used to reduce Euler discretization bias in early-exercise pricing.
      \end{block}
    \end{column}
    \begin{column}{0.31\textwidth}
      \begin{block}{Paper B}
        \paperlabel{AEMS / AEMS-SOR}\\
        Gatheral DMR. AES-type sampling is combined with a Milstein correction for a model with stochastic long-run variance.
      \end{block}
    \end{column}
    \begin{column}{0.31\textwidth}
      \begin{block}{Paper C}
        \paperlabel{Hybrid LSMC-PDE}\\
        GDMR Bermudan pricing. The stock-price direction is handled by a conditional FFT/PDE step and regression is performed over variance variables.
      \end{block}
    \end{column}
  \end{columns}
\end{frame}

\begin{frame}{Bermudan and American Pricing under Stochastic Volatility}
  \begin{columns}[T,onlytextwidth]
    \begin{column}{0.48\textwidth}
      \begin{block}{Pricing problem}
        For a put payoff \(h(S_t)=(K-S_t)^+\), Bermudan/American pricing requires
        \[
          V_{t_k}=\max\{h(S_{t_k}), C(t_k,Y_{t_k})\},
        \]
        where
        \[
          C(t_k,Y_{t_k})
          =e^{-r\dt}\E^\Q\!\left[V_{t_{k+1}}\mid Y_{t_k}\right].
        \]
      \end{block}
    \end{column}
    \begin{column}{0.48\textwidth}
      \begin{block}{Computational challenge}
        \begin{itemize}
          \item State variables include asset price and one or more variance factors.
          \item Simulation bias affects exercise decisions.
          \item Continuation values are conditional expectations, not explicit formulas.
        \end{itemize}
      \end{block}
    \end{column}
  \end{columns}

  \vspace{0.5em}
  \keymessage{The numerical price is driven by two decisions: how accurately paths are simulated, and how accurately continuation values are approximated.}
\end{frame}

\begin{frame}{Two Numerical Difficulties}
  \begin{columns}[T,onlytextwidth]
    \begin{column}{0.47\textwidth}
      \begin{block}{1. Variance simulation}
        Euler-type schemes can create negative variance approximations and introduce discretization bias.
        \[
          \dd v_t=\kappa(\theta-v_t)\dd t+\xi\sqrt{v_t}\dd W_t.
        \]
        For CIR-type variance factors, the next variance value can be sampled from its known noncentral chi-square distribution.
      \end{block}
    \end{column}
    \begin{column}{0.47\textwidth}
      \begin{block}{2. Continuation regression}
        In multifactor stochastic volatility,
        {\small
        \[
          C(t,Y_t)=\E^\Q[\hbox{discounted future value}\mid Y_t].
        \]}
        Plain LSMC must approximate this high-dimensional object from simulated paths.
      \end{block}
    \end{column}
  \end{columns}

  \vspace{0.5em}
  \centering
  \small Papers A and B focus on the simulation of variance paths. Paper C focuses on the approximation of continuation values.
\end{frame}

\begin{frame}{Model Hierarchy}
  \begin{columns}[T,onlytextwidth]
    \begin{column}{0.31\textwidth}
      \begin{block}{Heston}
        {\scriptsize
        \[
          \dd S_t=rS_t\dd t+\sqrt{v_t}S_t\dd W_t^{(1)}
        \]
        \[
          \dd v_t=\kappa(\theta-v_t)\dd t+\xi\sqrt{v_t}\dd W_t^{(2)}
        \]}
      \end{block}
    \end{column}
    \begin{column}{0.31\textwidth}
      \begin{block}{Double Heston}
        {\tiny
        \[
          \dd S_t=rS_t\dd t+
          S_t\sum_{i=1}^2\sqrt{v_t^{(i)}}\dd W_t^{(i)}
        \]
        \[
        \begin{aligned}
          \dd v_t^{(i)}
          &=\kappa_i(\bar v_i-v_t^{(i)})\dd t\\
          &\quad+\gamma_i\sqrt{v_t^{(i)}}\dd W_t^{(i+2)},
          \quad i=1,2.
        \end{aligned}
        \]}
      \end{block}
    \end{column}
    \begin{column}{0.31\textwidth}
      \begin{block}{Gatheral DMR/GDMR}
        {\tiny
        \[
          \dd S_t=rS_t\dd t+S_t\sqrt{v_t}\dd W_t^{(1)}
        \]
        \[
          \dd v_t=\kappa_1(v'_t-v_t)\dd t+
          \xi_1 v_t^{\delta_1}\dd W_t^{(2)}
        \]
        \[
          \dd v'_t=\kappa_2(\theta-v'_t)\dd t+
          \xi_2(v'_t)^{\delta_2}\dd W_t^{(3)}
        \]}
      \end{block}
    \end{column}
  \end{columns}

  \vspace{0.35em}
  \begin{center}
    \small Paper B uses the square-root DMR case, \(\delta_1=\delta_2=\tfrac12\). Paper C treats the CEV-type GDMR extension.
  \end{center}
\end{frame}
